../cictikz/src/cictikz/data/tex/cictikz_preamble.tex

% A reference built on the difference between two threshold voltages.
% Included so it can be recognised, not so it can be copied.
%
% Deliberately drawn as the same loop as l3_gmcell, because that is the point
% of the slide: identical topology, and the only difference is where the number
% comes from. The GM cell gets it from a 4:1 size ratio, which lithography
% controls. This gets it from two different channel implants, which nothing
% controls together.
%
% MN1 (right) is a standard Vt device with its source at ground, diode
% connected. MN2 (left) is a native device with R in its source. The PMOS
% mirror forces equal currents, and with equal W/L the overdrives cancel, so
% I*R = V_t1 - V_t2.
%
% Row spacing matches \grid = 1.6 so a device drawn by the \lv*mos macros spans
% exactly one row.

\begin{circuitikz}[american, thick, transform shape, circuit ee IEC]

  ../cictikz/src/cictikz/data/tex/cictikz_lib.tex

  \def\xl{0}          % native device, R in its source
  \def\xr{\grid*2.5}  % standard Vt device, source at ground

  \def\ybot{2.0}      % NMOS source level     (top of R)
  \def\ybotd{3.6}     % NMOS drain level      (\ybot + \grid)
  \def\ytop{4.4}      % PMOS drain level
  \def\ytops{6.0}     % PMOS source level     (\ytop + \grid)

  % ---- Left branch: R, native NMOS, PMOS ----
  \draw (\xl,0) \vground;
  \draw (\xl,0) -- ++(0,0.4);
  \draw (\xl,0.4) \vresistor{$R$};
  \draw (resEnd) -- (\xl,\ybot);

  \draw (\xl,\ybot) \lvmnmos{MN2}{};
  \draw (\xl,\ybotd) -- (\xl,\ytop);
  \draw (\xl,\ytop) \lvmpmos{MP2}{};
  \draw (\xl,\ytops) \vsupply;

  % ---- Right branch: standard NMOS straight to ground, PMOS ----
  \draw (\xr,0) \vground;
  \draw (\xr,0) -- (\xr,\ybot);

  \draw (\xr,\ybot) \lvnmos{MN1}{};
  \draw (\xr,\ybotd) -- (\xr,\ytop);
  \draw (\xr,\ytop) \lvpmos{MP1}{};
  \draw (\xr,\ytops) \vsupply;

  % ---- Gate lines: NMOS diode connected on the right, PMOS on the left ----
  \draw (MN2.gate) -- (MN1.gate);
  \coordinate (gBot) at ($(MN2.gate)!0.5!(MN1.gate)$);
  \draw (gBot) |- (\xr,\ybotd);

  \draw (MP2.gate) -- (MP1.gate);
  \coordinate (gTop) at ($(MP2.gate)!0.5!(MP1.gate)$);
  \draw (gTop) |- (\xl,\ytop);

  % ---- What sets the current ----
  \node[anchor=south] at ($(gTop)+(0,0.15)$) {$1:1$};
  \node[anchor=north] at ($(gBot)+(0,-0.15)$) {same $W/L$};

  \node[anchor=east] at (\xl-0.45,\ybot+0.8) {native, $V_{t2}$};
  \node[anchor=west] at (\xr+0.45,\ybot+0.8) {standard, $V_{t1}$};

  % ---- The whole reference is this one voltage ----
  % Drawn as plain nodes rather than a circuitikz `v=' annotation, same as the
  % V_o marks in l3_gmcell.
  \node[anchor=east] at (\xl-0.55,1.9) {$+$};
  \node[anchor=east] at (\xl-0.55,1.1) {$\Delta V_{t}$};
  \node[anchor=east] at (\xl-0.55,0.3) {$-$};

\end{circuitikz}
\end{document}
